\section{Negative offsets on the primary scale}
\label{sec:primary-regimes}

On the primary scale the cubic term remains in the limiting boundary
drift. Its contribution gives opposite signs to the curvature
coefficient \(wK''-K'\) on the two branches. When the normalized
constant coefficient is bounded away from zero, we combine this
curvature estimate on bounded radii with contraction at early heights
and large exits. In particular,
\cref{lem:curvature-criterion} gives concavity of \(\mathcal L\)
through a root-free passage and up to simple or double boundary roots, and
\cref{lem:direct-adaptive-positive} treats the positive half-orbits that
leave this scale. When the constant coefficient is small compared with
\(q_1^{3/2}\), a boundary root approaches the central line, and
\cref{prop:weak-root-passage} applies on the linear scale.
